Table \ref{tab:comparison} consolidates the dataset-wide results of Sections \ref{sec:baseline} and \ref{sec:optimized}. The optimized pipeline improves mean Sensitivity by 132\% (from 0.107 to 0.248) and mean \gls{fom} by 46\% (from 0.193 to 0.283), while mean Accuracy, already saturated by class imbalance, improves only marginally (from 0.974 to 0.980).

\begin{table}[htbp]
\centering
\caption{Baseline vs. optimized pipeline, mean values over 484 patients.}
\label{tab:comparison}
\begin{tabular}{lccc}
\toprule
Metric & Baseline & Optimized & Change \\
\midrule
Sensitivity (Pcd) & $0.107$ & $0.248$ & $+132\%$ \\
\gls{fom} & $0.193$ & $0.283$ & $+46\%$ \\
Accuracy & $0.974$ & $0.980$ & $+0.7\%$ \\
\bottomrule
\end{tabular}
\end{table}

These gains match the two causes found in Section \ref{sec:baseline}. Switching to \gls{flair} removes the \gls{csf} hyperintensity that competed with the tumor for the brightest \gls{fcm} cluster, and dynamic slice selection reduces, without eliminating, the fraction of patients whose processed slice contains no tumor at all, from 8.1\% to 5.8\%. Both changes act on the input to the pipeline, not on the clustering or edge-extraction logic. This suggests that the original bottleneck was mainly about the input given to the pipeline -- the contrast mechanism and the slice chosen -- rather than the segmentation algorithm itself.

Accuracy stays above 0.97 in both pipelines even when baseline Sensitivity is as low as 11\%. This is a known limitation of pixel-wise accuracy on heavily imbalanced segmentation tasks: boundary pixels are a small fraction of the image, so a prediction that almost completely misses the tumor is still "correct" on most background pixels. This is why Sensitivity and \gls{fom}, not Accuracy, are the metrics that matter for clinical reliability, and why the baseline pipeline can look visually reasonable in Section \ref{sec:baseline} while still being clinically unusable.

The two pipelines also share a blind spot. The same worst-case patient identified in Sections \ref{sec:baseline} and \ref{sec:optimized} gets $\mathrm{FOM}=0$ under both the fixed central slice and the \gls{flair}-hyperintensity criterion, because neither heuristic lands on an axial slice that intersects the annotated lesion. This shows that dynamic slice selection, although more robust on average, is still a 2-D, volume-statistics-based proxy and not an explicit localization of the tumor; it gives no guarantee for small or eccentrically located lesions that do not stand out in the global hyperintensity histogram. The \gls{fcm} cluster-selection heuristic of Section \ref{sec:baseline} has a similar limitation: even with the 40\% area cap and the solidity-area score, it is still an intensity-and-shape-only criterion with no spatial or anatomical prior. The multimodal contrast change, not a change to this heuristic, is responsible for most of the measured improvement. Finally, the dilated evaluation metric of Section \ref{sec:optimized} makes Sensitivity easier to interpret clinically, but by construction it is more permissive than the strict metric of Section \ref{sec:baseline}. The \gls{fom}, computed identically on strict boundaries in both pipelines, is therefore the more conservative and directly comparable indicator of the 46\% improvement reported above. The optimized pipeline reaches these results at a similar computational cost (Section \ref{sec:optimized}), which matters given its intended use on large-scale clinical datasets.